\cleardoublepage % memastikan bab baru dimulai di halaman ganjil (kanan)
\chapter{PENDAHULUAN}
\label{chap:pendahuluan}
% --- Latar Belakang ---
\section{Latar Belakang}
Perkembangan teknologi digital telah membawa perubahan fundamental dalam cara negara membangun, mengelola, dan mempertahankan sistem pertahanannya.
Lingkungan strategis global saat ini ditandai oleh meningkatnya kompleksitas ancaman yang bersifat multidimensi, tidak hanya dalam bentuk ancaman militer konvensional, tetapi juga ancaman siber, disrupsi informasi, serta perang berbasis data dan teknologi.
Dalam konteks tersebut, keunggulan pertahanan tidak lagi semata ditentukan oleh kekuatan alutsista atau jumlah personel, melainkan oleh kemampuan mengelola informasi, mengintegrasikan sistem, serta mengambil keputusan secara cepat dan akurat berbasis data.

Transformasi digital menjadi agenda strategis bagi banyak negara dalam memperkuat kapabilitas pertahanan modern.
Negara-negara maju telah menjadikan teknologi digital, data, dan sistem terintegrasi sebagai fondasi utama dalam membangun \textit{smart defence}, \textit{network-centric warfare}, dan \textit{joint all-domain operations}.
Pendekatan ini menempatkan data sebagai aset strategis dan menjadikan interoperabilitas sistem lintas domain sebagai kunci keunggulan operasional.
Dengan demikian, transformasi digital di sektor pertahanan bukan sekadar proses digitalisasi administrasi, melainkan perubahan menyeluruh pada tata kelola, proses bisnis, arsitektur sistem, dan budaya organisasi pertahanan.

Di tingkat nasional, Pemerintah Republik Indonesia telah menetapkan transformasi digital sebagai agenda pembangunan jangka panjang melalui berbagai regulasi dan kebijakan strategis, seperti Sistem Pemerintahan Berbasis Elektronik (SPBE), Percepatan Transformasi Digital Nasional, serta Rencana Pembangunan Jangka Panjang Nasional (RPJPN) 2025–2045.
Kebijakan tersebut menegaskan bahwa seluruh kementerian dan lembaga, termasuk Kementerian Pertahanan Republik Indonesia (Kemhan RI), dituntut untuk membangun ekosistem digital yang terintegrasi, aman, dan berkelanjutan guna mendukung penyelenggaraan pemerintahan dan pembangunan nasional, termasuk dalam sektor pertahanan.

Namun demikian, hingga saat ini implementasi transformasi digital di lingkungan Kemhan RI masih menghadapi berbagai tantangan mendasar.
Sistem informasi dan aplikasi pertahanan masih berkembang secara terpisah di berbagai unit kerja dan matra, sehingga menimbulkan fragmentasi sistem, redundansi aplikasi, serta lemahnya integrasi dan interoperabilitas data.
Kondisi ini berdampak pada terbatasnya kemampuan analisis lintas domain, lambatnya proses pengambilan keputusan strategis, serta belum optimalnya pemanfaatan data sebagai aset pertahanan.
Selain itu, meningkatnya ancaman siber terhadap infrastruktur dan data pertahanan menuntut adanya fondasi arsitektur digital yang lebih kuat, terstandar, dan berorientasi pada keamanan sejak tahap perancangan.

Di sisi lain, Kemhan RI juga dihadapkan pada tuntutan untuk mewujudkan modernisasi pertahanan yang selaras dengan konsep pertahanan cerdas (\textit{smart defence}).
Konsep ini menekankan pentingnya integrasi antara manusia, proses, teknologi, dan kebijakan dalam satu arsitektur pertahanan digital yang adaptif dan resilien.
Tanpa adanya kerangka arsitektur transformasi digital yang jelas dan terstruktur, upaya modernisasi pertahanan berisiko berjalan parsial, tidak terarah, dan sulit diimplementasikan secara berkelanjutan.

Berdasarkan kondisi tersebut, diperlukan suatu pendekatan sistematis untuk merancang arsitektur transformasi digital yang dapat menjadi arah strategis bagi Kemhan RI.
Arsitektur ini diharapkan mampu mengintegrasikan aspek bisnis, data, aplikasi, dan teknologi pertahanan secara holistik, sekaligus selaras dengan kebijakan nasional dan praktik terbaik internasional.
Penggunaan \textit{enterprise architecture framework} seperti TOGAF menjadi relevan karena menyediakan metodologi terstruktur untuk menyelaraskan kebutuhan strategis organisasi dengan solusi teknologi secara menyeluruh.

Oleh karena itu, penelitian ini difokuskan pada perancangan arsitektur transformasi digital pada Kementerian Pertahanan Republik Indonesia sebagai landasan dalam mewujudkan pertahanan cerdas yang modern, terintegrasi, dan berdaya tahan tinggi terhadap dinamika ancaman di era digital.
Hasil penelitian diharapkan dapat memberikan kontribusi strategis, baik secara akademik maupun praktis, dalam mendukung pembangunan ekosistem digital pertahanan nasional yang berkelanjutan.

% --- Rumusan Masalah ---
\section{Rumusan Masalah}
Berdasarkan latar belakang yang telah diuraikan, terdapat sejumlah permasalahan mendasar yang menghambat percepatan transformasi digital di lingkungan Kementerian Pertahanan Republik Indonesia.
Permasalahan tersebut muncul sebagai konsekuensi dari fragmentasi sistem, belum optimalnya tata kelola data dan teknologi, serta belum adanya kerangka arsitektur transformasi digital yang dapat menjadi acuan terintegrasi.
Sampai dengan saat ini sistem informasi pertahanan masih berjalan secara terpisah antar unit dan matra, sehingga menghambat integrasi data dan kolaborasi lintas organisasi.
Selain itu, kemajuan teknologi dan meningkatnya ancaman siber menuntut Kemhan untuk membangun fondasi digital yang lebih kuat, terarah, dan sesuai dengan mandat regulatif nasional.

Untuk menjawab tantangan tersebut, penelitian ini merumuskan beberapa pertanyaan pokok sebagai berikut:
\begin{enumerate}
\item	Bagaimana kondisi eksisting (as-is) proses bisnis, data, aplikasi dan teknologi di lingkungan Kementerian Pertahanan yang berkaitan dengan transformasi digital?
\item	Bagaimana merancang arsitektur transformasi digital Kemhan RI yang relevan dengan kebutuhan pertahanan modern?
\end{enumerate}
Rumusan masalah tersebut menjadi dasar bagi penelitian ini dalam menghasilkan arsitektur transformasi digital yang dapat digunakan Kemhan RI sebagai pedoman strategis dalam membangun ekosistem digital pertahanan yang modern, aman, dan berkelanjutan.

% --- Tujuan ---
\section{Tujuan Penelitian}
Penelitian ini bertujuan untuk merancang suatu arsitektur transformasi digital yang komprehensif bagi Kementerian Pertahanan Republik Indonesia sebagai landasan strategis dalam memperkuat tata kelola pertahanan modern.
Upaya ini dilakukan untuk menjawab tantangan fragmentasi sistem, lemahnya integrasi data pertahanan, meningkatnya ancaman siber, serta tuntutan regulatif nasional yang mewajibkan Kemhan memiliki arsitektur dan peta rencana SPBE yang terstruktur dan terintegrasi.
Secara khusus, penelitian ini memiliki tujuan sebagai berikut:
\begin{enumerate}
    \item Mengidentifikasi kondisi eksisting (as-is) ekosistem digital di lingkungan Kementerian Pertahanan, mencakup proses bisnis, aliran data, aplikasi, infrastruktur teknologi 
    terhadap transformasi digital.
    \item Merancang arsitektur transformasi digital Kementerian Pertahanan berdasarkan kerangka TOGAF yang mencakup empat domain inti,
    yaitu: arsitektur bisnis, data, aplikasi, dan teknologi.
    \item Memberikan rekomendasi strategis berupa roadmap implementasi transformasi digital Kemhan RI yang berfungsi sebagai
    panduan jangka menengah bagi organisasi dalam mengembangkan dan mengimplementasikan berbagai inisiatif digital secara bertahap, terarah, dan berkelanjutan.
\end{enumerate}
Melalui pencapaian tujuan-tujuan tersebut, penelitian ini diharapkan dapat memberikan kontribusi nyata dalam mendukung modernisasi pertahanan nasional
melalui pembangunan arsitektur digital yang kuat, adaptif, dan selaras dengan kebutuhan keamanan nasional di era digital.

% --- Manfaat Penelitian ---
\section{Manfaat Penelitian}
Penelitian mengenai desain arsitektur transformasi digital pada Kementerian Pertahanan RI diharapkan memberikan kontribusi teoretis 
maupun praktis yang relevan bagi pengembangan ekosistem digital pertahanan nasional.

\subsection{Manfaat Teoritis}
\begin{enumerate}
    \item Menghasilkan desain arsitektur transformasi digital pada Kementerian Pertahanan dengan menggunakan kerangka kerja TOGAF.
    \item Menjadi referensi akademik bagi penelitian lanjutan dalam bidang \textit{smart defense}, \textit{AI governance}, dan \textit{digital military transformation}.
\end{enumerate}

\subsection{Manfaat Praktis}
\begin{enumerate}
    \item Memberikan rekomendasi dan model desain arsitektur transformasi digital bagi Kementerian Pertahanan.
    \item Mendukung peningkatan tata kelola organisasi, integrasi sistem pertahanan dan modernisasi pertahanan.
    \item Memperkuat kapabilitas keamanan siber Kemhan.
\end{enumerate}
Dengan demikian, penelitian ini tidak hanya memberikan nilai ilmiah, tetapi juga memberikan arah strategis yang dapat langsung diterapkan
pada upaya transformasi digital di lingkungan Kementerian Pertahanan RI.

% --- Ruang Lingkup Penelitian ---
\section{Ruang Lingkup Penelitian}
Ruang lingkup penelitian ini ditetapkan untuk memastikan fokus analisis yang jelas dan sesuai dengan tujuan penelitian.
Adapun batasan dan cakupan penelitian ini meliputi:
\begin{enumerate}
    \item Lingkup organisasi: Kementerian Pertahanan Republik Indonesia, khususnya unit-unit kerja yang terkait dengan perencanaan,
    pengelolaan dan implementasi transformasi digital, seperti Ditjen Pothan, Pusdatin Kemhan dan Bainstrahan Kemhan.
    \item Lingkup Teknis: Arsitektur bisnis, data, aplikasi, dan teknologi berdasarkan kerangka kerja TOGAF ADM.
    \item Batasan: Penelitian ini tidak mencakup aspek rahasia militer dan operasi tempur, namun fokus pada integrasi sistem informasi, keamanan data, dan tata kelola digital.
\end{enumerate}

% --- Hipotesis ---
\section{Hipotesis}
Hipotesis penelitian diidentifikasi dari latar belakang, rumusan masalah dan tujuan penelitian.

\textbf{Premis 1:} Transformasi digital pada sektor pertahanan memerlukan integrasi antara proses bisnis, data, aplikasi, teknologi, tata kelola, dan keamanan siber untuk mendukung interoperabilitas sistem, pengambilan keputusan berbasis data, serta peningkatan efektivitas organisasi.
Berbagai penelitian menunjukkan bahwa \textit{Enterprise Architecture} (EA) mampu menjadi instrumen strategis dalam menyelaraskan kebutuhan organisasi dengan pemanfaatan teknologi digital secara terstruktur dan berkelanjutan \autocite{Syamfithriani,Mattila2022}.

\textbf{Premis 2:} Kementerian Pertahanan Republik Indonesia masih menghadapi tantangan berupa fragmentasi sistem informasi, silo data antar unit kerja, rendahnya interoperabilitas sistem, serta belum tersedianya kerangka arsitektur transformasi digital yang terintegrasi dan selaras dengan kebijakan SPBE Nasional.
Kondisi tersebut menunjukkan perlunya suatu \textit{framework} arsitektur yang mampu mengintegrasikan aspek bisnis, data, aplikasi, dan teknologi dalam mendukung transformasi digital pertahanan.

\textbf{Premis 3:} TOGAF ADM standar dirancang sebagai kerangka \textit{Enterprise Architecture} yang bersifat umum dan belum secara eksplisit mengakomodasi kebutuhan operasional pertahanan seperti C5ISR, operasi multi-domain, interoperabilitas lintas matra, dan prinsip \textit{security-by-design}.
Oleh karena itu diperlukan modifikasi \textit{framework} agar lebih sesuai dengan karakteristik dan kebutuhan transformasi digital sektor pertahanan Indonesia.

Berdasarkan premis-premis tersebut, diusulkan hipotesis penelitian ini sebagai berikut:
\begin{enumerate}
    \item \textbf{Hipotesis 1:} \textit{Framework} TOGAF Pertahanan Indonesia yang dirancang berdasarkan pendekatan TOGAF ADM mampu menyediakan arsitektur transformasi digital yang lebih terintegrasi, selaras dengan kebutuhan organisasi, serta layak digunakan sebagai \textit{blueprint} transformasi digital Kementerian Pertahanan Republik Indonesia.
    \item \textbf{Hipotesis 2:} Penambahan fase C5ISR \textit{Architecture} pada TOGAF Pertahanan Indonesia akan menghasilkan arsitektur yang lebih relevan dalam mendukung interoperabilitas lintas unit dan matra, keamanan siber, integrasi data pertahanan, serta kebutuhan operasi pertahanan modern dibandingkan TOGAF ADM standar.
\end{enumerate}

% % --- Keluaran Penelitian ---
% \section{Keluaran Penelitian}
% Penelitian ini diarahkan untuk menghasilkan keluaran yang bersifat konseptual maupun aplikatif dalam mendukung transformasi digital Kementerian Pertahanan RI.
% Adapun keluaran yang dihasilkan meliputi:
% \begin{enumerate}
%     \item Model arsitektur transformasi digital Kemhan RI dengan menggunakan pendekatan TOGAF ADM Phase A–D, mencakup arsitektur bisnis, data, aplikasi, dan teknologi.
%     \item Rekomendasi Strategis dan Roadmap Implementasi Transformasi Digital Kemhan 2026–2030.
% \end{enumerate}


% --- Metodologi Penelitian ---
\section{Metodologi Penelitian}

\subsection{Pendekatan Penelitian}
Pendekatan penelitian yang digunakan dalam studi ini adalah \textit{Design Science Research Methodology} (DSRM).
DSRM merupakan pendekatan penelitian yang berorientasi pada \textit{problem-solving} melalui pengembangan artefak yang dapat memberikan solusi terhadap permasalahan nyata dalam suatu organisasi.
Pendekatan ini sangat relevan digunakan dalam penelitian yang bertujuan merancang model, kerangka kerja, atau arsitektur yang bersifat implementatif.

Menurut \textcite{Hevner2004}, \textit{Design Science} bertujuan menghasilkan artefak yang efektif dan dapat diuji, seperti model, metode, ataupun arsitektur.
Pendekatan ini berupaya menggabungkan proses ilmiah dengan kebutuhan praktis organisasi, sehingga artefak yang dihasilkan tidak hanya memiliki nilai teoritis, tetapi juga memberikan manfaat aplikatif.
Dalam konteks penelitian ini, artefak yang dikembangkan berupa arsitektur transformasi digital Kemhan RI yang disusun secara terstruktur dengan kerangka TOGAF ADM.

DSRM terdiri dari enam tahapan utama, yaitu: (1) identifikasi masalah dan motivasi, (2) penetapan tujuan solusi, (3) perancangan dan pengembangan artefak, (4) demonstrasi terhadap konteks atau skenario penggunaan, (5) evaluasi artefak, dan (6) komunikasi hasil penelitian.
Keenam tahap ini memberikan alur kerja sistematis yang memandu peneliti sejak perumusan masalah hingga penyebaran hasil.

Pendekatan DSRM dipilih karena beberapa alasan utama:
\begin{enumerate}
  \item Penelitian ini berfokus pada pengembangan arsitektur transformasi digital pada Kementerian Pertahanan yang merupakan artefak desain.
  DSRM menyediakan kerangka yang tepat untuk merancang artefak dengan mempertimbangkan kebutuhan organisasi dan teori yang mendukung.
  \item Transformasi digital Kemhan RI membutuhkan pendekatan yang tidak hanya teoritis, tetapi juga aplikatif.
  DSRM memastikan bahwa artefak yang dihasilkan dapat diterapkan pada konteks nyata organisasi pertahanan.
  \item Melalui tahapan evaluasi, artefak yang dikembangkan dapat dinilai kesesuaiannya oleh para pakar, unit terkait, dan pemangku kepentingan di Kemhan RI.
  Dengan demikian, model yang dihasilkan memiliki validitas yang kuat untuk digunakan dalam desain arsitektur transformasi digital Kementerian Pertahanan.
\end{enumerate}
Dengan demikian, penggunaan DSRM sebagai pendekatan penelitian memberikan landasan metodologis yang kuat dan sesuai untuk menghasilkan arsitektur transformasi digital Kementerian Pertahanan yang sistematis dan komprehensif.

\subsection{Tahapan Penelitian}
Penelitian ini menggunakan \textit{Design Science Research Methodology} (DSRM) yang terdiri dari enam tahapan utama sebagaimana dikemukakan oleh Peffers dkk. (2007).
Setiap tahapan dalam DSRM memberikan struktur penelitian yang sistematis dan berorientasi pada pengembangan artefak yang dapat menyelesaikan permasalahan nyata, dalam hal ini arsitektur transformasi digital Kemhan RI.
Adapun tahapan penelitian tersebut dijelaskan sebagai berikut:

\subsubsection{Tahap 1 – Problem Identification and Motivation}
Tahap pertama bertujuan untuk mengidentifikasi permasalahan inti yang menghambat transformasi digital di lingkungan Kemhan RI serta menentukan urgensi penelitian.
Berdasarkan analisis dokumen dan kebijakan, permasalahan utama yang ditemukan adalah sebagai berikut:
\begin{enumerate}
  \item Fragmentasi sistem dan data pertahanan.
  Sistem informasi di Kemhan dan angkatan berjalan sendiri-sendiri, menyebabkan rendahnya interoperabilitas dan integrasi data.
  \item Belum adanya \textit{blueprint} arsitektur digital Kemhan.
  Kemhan belum memiliki arsitektur SPBE dan arsitektur pertahanan yang terstruktur untuk mendukung modernisasi digital.
  \item Ancaman siber yang meningkat.
  \textcite{KemhanRI2014} menunjukkan bahwa ancaman siber menjadi tantangan strategis yang belum diimbangi dengan kesiapan sistem dan SDM dalam pertahanan siber.
  \item Keterbatasan SDM digital dan tata kelola TI.
  Kompetensi digital belum merata, dan belum ada mekanisme tata kelola arsitektur yang menyeluruh.
  \item Tuntutan regulasi nasional.
  \textcite{KementerianKomunikasidanInformatika2018}, \textcite{RepublikIndonesia2023}, \textcite{RI2019}, dan \textcite{RPJPN2025} mewajibkan Kemhan membangun ekosistem digital pertahanan yang terintegrasi.
\end{enumerate}

\subsubsection{Tahap 2 – Defining the Objectives for a Solution}
Pada tahap ini ditetapkan tujuan solusi yang ingin dicapai oleh artefak arsitektur.
Tujuan tersebut disusun berdasarkan permasalahan yang telah diidentifikasi, kebutuhan organisasi, serta mandat regulasi. Adapun tujuan solusi adalah:
\begin{enumerate}
  \item Menghasilkan arsitektur bisnis, data, aplikasi, dan teknologi yang terintegrasi berdasarkan kerangka TOGAF.
  \item Menyediakan \textit{blueprint} yang dapat digunakan sebagai pedoman transformasi digital Kemhan RI.
  \item Memperkuat tata kelola data pertahanan dan meningkatkan interoperabilitas antar unit organisasi dan matra.
  \item Membangun fondasi penguatan keamanan siber Kemhan.
  \item Mendukung penyusunan \textit{roadmap} implementasi transformasi digital Kemhan RI 2026–2030.
\end{enumerate}
Tujuan-tujuan ini menjadi dasar dalam pengembangan artefak dan memastikan bahwa solusi yang dihasilkan relevan dengan kondisi aktual dan kebutuhan jangka panjang organisasi.

\subsubsection{Tahap 3 – Design and Development}
Pada tahap ini dilakukan perancangan dan pengembangan artefak arsitektur transformasi digital Kemhan RI menggunakan tahapan TOGAF ADM.
Ringkasan perancangan arsitektur, mencakup tujuan utama, aktivitas kunci, dan hasil akhir pada setiap tahapan TOGAF, disajikan pada Tabel \ref{tab:perancangan-arsitektur}.

\begin{longtable}{|c|p{2.4cm}|p{3cm}|p{3.6cm}|p{2.8cm}|}
\caption{Perancangan Arsitektur Transformasi Digital Kemhan RI}
\label{tab:perancangan-arsitektur} \\
\hline
\textbf{No.} & \textbf{Tahapan TOGAF} & \textbf{Tujuan Utama} & \textbf{Aktivitas Kunci} & \textbf{Hasil Akhir} \\
\hline
\endfirsthead

\hline
\textbf{No.} & \textbf{Tahapan TOGAF} & \textbf{Tujuan Utama} & \textbf{Aktivitas Kunci} & \textbf{Hasil Akhir} \\
\hline
\endhead

1 & \textit{Preliminary Phase} & Menentukan ruang lingkup, prinsip arsitektur, dan kesiapan organisasi dalam pengembangan arsitektur transformasi digital.
& Identifikasi ruang lingkup transformasi digital pertahanan; identifikasi keterkaitan dengan regulasi nasional.
& Dokumen hasil identifikasi ruang lingkup dan regulasi nasional. \\
\hline
2 & \textit{Architecture Vision} & Merumuskan visi, misi, dan tujuan strategis arsitektur untuk mendukung transformasi digital pertahanan.
& Menyusun visi dan misi arsitektur pertahanan digital; mengidentifikasi kebutuhan strategis organisasi; identifikasi pemangku kepentingan utama.
& Visi \& Misi Pertahanan Digital. \\
\hline
3 & \textit{Business Architecture} & Merancang model proses bisnis pertahanan yang efektif untuk mendukung komando, kendali, integrasi data, serta kolaborasi lintas satuan.
& Pemodelan proses bisnis inti dan pendukung; analisis kesenjangan proses (\textit{as-is} vs. \textit{to-be}); identifikasi domain layanan digital pertahanan; penyusunan \textit{business process model} ekosistem digital Kemhan.
& \textit{Capability Mapping} Kementerian Pertahanan. \\
\hline
4 & \textit{Information Systems Architecture} & Menentukan kebutuhan arsitektur data dan aplikasi yang mendukung proses bisnis pertahanan.
& Pemodelan entitas data utama dan hubungan antar sistem; desain aplikasi pendukung operasi pertahanan digital.
& Model arsitektur data dan aplikasi pertahanan (\textit{Data Model, Application Portfolio}). \\
\hline
5 & \textit{Technology Architecture} & Menyusun arsitektur teknologi yang mencakup infrastruktur, platform, dan keamanan siber sebagai fondasi ekosistem digital pertahanan.
& Mendesain arsitektur teknologi target (jaringan, server, \textit{cloud}, keamanan siber); menentukan standar teknologi dan protokol keamanan.
& Arsitektur teknologi pertahanan (\textit{technology architecture model}). \\
\hline
\end{longtable}

\subsubsection{Tahap 4 – Demonstration}
Pada tahap ini, artefak yang telah dikembangkan diuji melalui proses demonstrasi untuk menunjukkan bagaimana artefak tersebut menyelesaikan permasalahan yang telah diidentifikasi. Demonstrasi dilakukan melalui:
\begin{enumerate}
  \item Diskusi kelompok terarah (FGD) dengan pakar dan pejabat Kemhan untuk melihat sejauh mana arsitektur yang dirancang dapat diimplementasikan secara praktis.
\end{enumerate}
Demonstrasi bertujuan memastikan bahwa artefak dapat berfungsi secara logis dalam konteks operasional Kemhan RI dan memberikan gambaran implementasi awal sebelum dilakukan evaluasi mendalam.

\subsubsection{Tahap 5 – Evaluation}
Tahapan ini bertujuan untuk menilai efektivitas artefak arsitektur transformasi digital Kemhan RI dalam menjawab permasalahan dan mencapai tujuan penelitian. Evaluasi dilakukan dengan dua pendekatan utama:
\begin{enumerate}
  \item Evaluasi Teoritis:
  \begin{enumerate}
    \item Menilai kelayakan teknis dan konseptual dari arsitektur transformasi digital yang dirancang, termasuk konsistensi dengan prinsip-prinsip arsitektur \textit{enterprise} (EA) dan keselarasan dengan kerangka kerja TOGAF ADM.
    \item Menilai kesesuaian artefak dengan standar dan kebijakan nasional, seperti arsitektur SPBE Nasional dan regulasi transformasi digital pemerintahan.
    \item Memastikan validitas ilmiah melalui peninjauan oleh pakar akademik dalam bidang sistem informasi, \textit{enterprise architecture}, dan transformasi digital.
  \end{enumerate}
  \item Evaluasi Praktis:
  \begin{enumerate}
    \item Melibatkan pakar dan praktisi dari lingkungan pertahanan, seperti STEI ITB, Ditjen Potensi Pertahanan, Pusdatin Kemhan dan Bainstrahan Kemhan yang relevan untuk memberikan penilaian aplikatif terhadap artefak.
    \item Penilaian dilakukan pada aspek-aspek implementatif, termasuk: efisiensi proses; interoperabilitas sistem pendukung; keamanan dan kesiapan pengelolaan data; serta kelayakan adopsi dalam struktur organisasi Kemhan.
    \item Mengidentifikasi area perbaikan agar arsitektur dapat disesuaikan dengan kondisi operasional Kemhan RI dan kebutuhan transformasi pertahanan.
  \end{enumerate}
\end{enumerate}
Hasil evaluasi memberikan dasar untuk melakukan revisi dan penyempurnaan pada model agar sesuai dengan kebutuhan organisasi dan dinamika lingkungan pertahanan.

\subsubsection{Tahap 6 – Communication}
Tahap terakhir adalah komunikasi hasil penelitian. Pada tahap ini, artefak dan temuan penelitian disajikan dalam bentuk laporan ilmiah, presentasi kepada pemangku kepentingan, serta publikasi akademik apabila relevan.
Komunikasi bertujuan memastikan bahwa hasil penelitian dapat dipahami, digunakan, dan diadopsi oleh pihak yang berkepentingan.
Dalam konteks penelitian ini, komunikasi melibatkan penyampaian model arsitektur transformasi digital Kemhan RI dan \textit{roadmap} implementasi kepada Kemhan RI sebagai pengguna utama.
Dengan demikian, artefak tidak hanya menjadi kontribusi teoretis, tetapi juga memberikan manfaat praktis bagi organisasi.

\subsection{Jenis dan Sumber Data}
Penelitian ini menggunakan dua jenis data utama, yaitu data primer dan data sekunder, yang dikumpulkan untuk mendukung proses analisis pada setiap tahapan DSRM.
Penggunaan kedua jenis data ini diperlukan agar perancangan arsitektur transformasi digital Kemhan RI memiliki landasan konseptual yang kuat sekaligus relevansi praktis terhadap kondisi aktual di Kementerian Pertahanan RI.
\begin{enumerate}
  \item Data Primer.
  Data primer merupakan data yang diperoleh langsung dari sumber pertama melalui interaksi dengan pemangku kepentingan.
  Data ini diperlukan untuk memahami secara mendalam kondisi aktual sistem informasi, proses bisnis, kebutuhan organisasi, serta tantangan transformasi digital di lingkungan Kemhan. Data primer meliputi:
  \begin{enumerate}
    \item Hasil wawancara dengan pejabat atau staf Kemhan yang memahami sistem informasi, teknologi pertahanan, atau kebijakan SPBE.
    \item Validasi pakar (\textit{expert judgement}) melalui akademisi atau ahli di bidang \textit{enterprise architecture} dan transformasi digital pertahanan.
  \end{enumerate}
  Data primer ini sangat penting untuk mengidentifikasi kesenjangan (\textit{gap}) antara kondisi eksisting dan kebutuhan arsitektur target.
  \item Data Sekunder.
  Data sekunder diperoleh dari dokumen dan sumber referensi yang relevan, antara lain:
  \begin{enumerate}
    \item Peraturan Presiden No. 95 Tahun 2018 tentang SPBE.
    \item Peraturan Presiden No. 39 Tahun 2019 tentang Satu Data Indonesia.
    \item Peraturan Presiden No. 82 Tahun 2023 tentang Percepatan Transformasi Digital.
    \item Rencana Pembangunan Jangka Panjang Nasional (RPJPN) 2025--2045.
    \item Peraturan Menteri Pertahanan No. 82 Tahun 2014 tentang Pertahanan Siber.
    \item Undang-Undang No. 23 Tahun 2019 tentang Pengelolaan Sumber Daya Nasional untuk Pertahanan Negara.
    \item Rencana Strategis Kementerian Pertahanan Tahun 2020--2024.
    \item Jurnal ilmiah, buku, dan publikasi akademik yang membahas transformasi digital, \textit{enterprise architecture}, kompetensi digital, serta pengembangan SDM di sektor pertahanan.
    \item Referensi internasional, termasuk \textit{DoD Digital Modernization Strategy} dan penelitian global mengenai \textit{digital defense capability}, untuk membandingkan praktik terbaik dalam transformasi digital pertahanan.
  \end{enumerate}
  Data sekunder digunakan sebagai landasan teoretis untuk menyusun kerangka konseptual, mendukung analisis \textit{gap} kompetensi, serta memastikan bahwa arsitektur yang dihasilkan selaras dengan perkembangan global dan kebijakan nasional.
\end{enumerate}

\subsection{Metode Pengumpulan Data}
Metode pengumpulan data dilakukan melalui beberapa tahapan:
\begin{enumerate}
  \item Studi Literatur.
  Mengkaji teori dasar EA, TOGAF, \textit{Smart Defense}, dan Transformasi Digital.
  \item Analisis Dokumen Internal.
  Analisis dokumen dilakukan terhadap data dan kebijakan yang dikeluarkan oleh Kemhan RI.
  \item Wawancara Terarah.
  Wawancara dilakukan dengan pejabat atau personel dari unit-unit Kemhan yang relevan, seperti Ditjen Potensi Pertahanan, Pusdatin Kemhan dan Bainstrahan Kemhan. Wawancara ini bertujuan untuk:
  \begin{enumerate}
    \item Menggali informasi mengenai tantangan implementasi transformasi digital.
    \item Mendapatkan perspektif mengenai arah transformasi digital Kemhan RI.
    \item Memvalidasi temuan awal dari analisis dokumen.
  \end{enumerate}
  Metode wawancara semi-terstruktur memungkinkan penggalian informasi yang fleksibel namun tetap terfokus pada isu penelitian.
  \item \textit{Focus Group Discussion} (FGD).
  FGD digunakan pada tahap evaluasi artefak, sebagai bagian dari pendekatan evaluasi praktis pada DSRM. FGD dilakukan dengan melibatkan akademisi, Ditjen Pothan, Pusdatin Kemhan, dan Bainstrahan Kemhan.
  FGD bertujuan menilai kelayakan arsitektur transformasi digital dari aspek teknis, organisasi, keamanan, dan kesiapan implementasi.
\end{enumerate}


% --- Sistematika Penulisan ---
\section{Sistematika Penulisan}
Sistematika penulisan dalam proposal penelitian ini disusun untuk memberikan alur pembahasan yang runtut, logis, dan mudah dipahami.
Setiap bab dirancang untuk saling mendukung satu sama lain, mulai dari pengenalan masalah hingga penyusunan metodologi penelitian yang digunakan dalam merancang arsitektur transformasi digital pada Kementerian Pertahanan RI.
Adapun sistematika penulisan proposal ini adalah sebagai berikut:
\begin{enumerate}
\item	Bab I Pendahuluan: Bab ini menjelaskan landasan awal penelitian, meliputi latar belakang permasalahan, rumusan masalah, tujuan penelitian, manfaat penelitian, ruang lingkup penelitian, metodologi penelitian, serta sistematika penulisan.
Bab ini juga menguraikan pendekatan \textit{Design Science Research Methodology} (DSRM) beserta tahapan-tahapannya, jenis dan sumber data, serta metode pengumpulan data yang digunakan.
\item	Bab II Tinjauan Pustaka: Bab ini menyajikan kajian teoretis dan konseptual yang menjadi dasar penelitian.
Pembahasan mencakup konsep transformasi digital, transformasi digital di sektor pertahanan, arah dan tantangan transformasi digital Kementerian Pertahanan RI,
serta konsep pertahanan cerdas (\textit{smart defence}). Selain itu, bab ini juga mengulas \textit{framework} TOGAF sebagai kerangka kerja arsitektur \textit{enterprise} yang digunakan dalam perancangan arsitektur transformasi digital.
Tinjauan pustaka ini berfungsi sebagai pijakan teoritis dan konseptual dalam pengembangan artefak penelitian.
\item	Bab III Analisis dan Desain: Bab ini menyajikan analisis kondisi eksisting Kemhan RI, mencakup analisis organisasi dan tata kelola, regulasi dan kebijakan, permasalahan transformasi digital, serta kebutuhan \textit{smart defense}.
Berdasarkan analisis tersebut, bab ini merancang kerangka TOGAF Pertahanan Indonesia beserta prinsip arsitektur dan keseluruhan fase pengembangannya sebagai artefak utama penelitian.
\item	Bab IV Rencana Evaluasi: Bab ini menjelaskan metode evaluasi artefak arsitektur yang dirancang, meliputi penilaian menggunakan \textit{Enterprise Architecture Scorecard}\texttrademark{} dan \textit{expert judgement} melalui FGD, beserta kriteria keberhasilan evaluasi.
\item   Daftar Pustaka: Berisi daftar referensi yang digunakan dalam penyusunan laporan tesis.
\end{enumerate}